% !TeX program = pdflatex
% ======================================================================
% The Relational Spin-Closure Hypothesis
% Mathematical paper draft, dependency-ordered and theorem-centred.
%
% Source of mathematical truth for this draft:
% current Lean source in
% https://github.com/antaris82/The_Relational_Spin-Closure_Hypothesis
%
% Editorial attribution:
% Jan Seeck -- Editor
% ======================================================================

\documentclass[11pt,a4paper]{article}

\usepackage[margin=27mm]{geometry}
\usepackage[T1]{fontenc}
\usepackage[utf8]{inputenc}
\usepackage{microtype}
\usepackage{amsmath,amssymb,amsthm,mathtools}
\usepackage{enumitem}
\usepackage{booktabs}
\usepackage{array}
\usepackage{xcolor}
\usepackage{url}
\usepackage[
  hidelinks,
  pdftitle={The Relational Spin-Closure Hypothesis: From Euclidean Three-Space to Conditional Lorentzian Spin Geometry},
  pdfauthor={Jan Seeck (Editor)},
  pdfsubject={Conditional reconstruction from Euclidean three-space to Lorentzian Spin geometry}
]{hyperref}
\usepackage[nameinlink,noabbrev]{cleveref}

\setlength{\parindent}{1.4em}
\setlength{\parskip}{0pt}
\setlist[itemize]{topsep=.4em,itemsep=.2em}
\setlist[enumerate]{topsep=.4em,itemsep=.2em}

% ----------------------------------------------------------------------
% Theorem environments
% ----------------------------------------------------------------------

\newtheorem{theorem}{Theorem}[section]
\newtheorem{proposition}[theorem]{Proposition}
\newtheorem{lemma}[theorem]{Lemma}
\newtheorem{corollary}[theorem]{Corollary}

\theoremstyle{definition}
\newtheorem{definition}[theorem]{Definition}
\newtheorem{construction}[theorem]{Construction}
\newtheorem{hypothesis}[theorem]{Hypothesis}

\theoremstyle{remark}
\newtheorem{remark}[theorem]{Remark}
\newtheorem{boundary}[theorem]{Boundary statement}
\newtheorem{control}[theorem]{Control}

% ----------------------------------------------------------------------
% Notation
% ----------------------------------------------------------------------

\newcommand{\R}{\mathbb{R}}
\newcommand{\Ztwo}{\mathbb{Z}/2}
\newcommand{\Smodel}{\mathcal{S}}
\newcommand{\Cl}{\operatorname{Cl}}
\newcommand{\SpinG}{G_{\mathrm{Spin}}}
\newcommand{\LorG}{G_{\mathrm{Lor}}}
\newcommand{\End}{\operatorname{End}}
\newcommand{\GL}{\operatorname{GL}}
\newcommand{\id}{\operatorname{id}}
\newcommand{\im}{\operatorname{im}}
\newcommand{\kerop}{\operatorname{ker}}
\newcommand{\Ad}{\operatorname{Ad}}
\newcommand{\Int}{\operatorname{Int}}
\newcommand{\Diff}{\operatorname{Diff}}
\newcommand{\Cech}{\ifmmode\text{\v{C}ech}\else\v{C}ech\fi}
\newcommand{\cU}{\mathcal{U}}
\newcommand{\cR}{\mathcal{R}}
\newcommand{\cE}{\mathcal{E}}
\newcommand{\gB}{\mathfrak{g}_{B}}
\newcommand{\wedgeS}{\Lambda^{2}\Smodel}
\newcommand{\dd}{\,\mathrm{d}}

\newcommand{\statusderived}{\textsc{derived}}
\newcommand{\statusdatum}{\textsc{additional datum}}
\newcommand{\statusopen}{\textsc{open}}
\newcommand{\statusinfra}{\textsc{infrastructure-limited}}

\DeclareMathOperator{\detop}{det}

% ----------------------------------------------------------------------
% Title
% ----------------------------------------------------------------------

\title{\textbf{The Relational Spin-Closure Hypothesis}\\[0.35em]
\large From Euclidean Three-Space to Conditional Lorentzian Spin Geometry}

\author{Edited by Jan Seeck}
\date{}

\begin{document}
\maketitle

\begin{abstract}
We study a reconstruction problem whose primitive local input is only the
standard positive-definite Euclidean space \(V=\R^{3}\).
From this datum one obtains the real spin factor
\(\Smodel=\R\oplus V\), its quadratic form
\(N(a,v)=a^{2}-\langle v,v\rangle\) of signature \((1,3)\), the real
Clifford algebra \(\Cl(V,q)\), an intrinsic Spin group
\(\SpinG\), and a surjective central double cover
\[
   1\longrightarrow\{\pm1\}
   \longrightarrow \SpinG
   \overset{\rho}{\longrightarrow}\LorG
   \longrightarrow 1
\]
of an intrinsically defined proper cone-preserving Lorentz group.
In parallel, the Lorentz-skew endomorphisms of \(\Smodel\) form a
six-dimensional Lie algebra \(\gB\) linearly equivalent to
\(\Lambda^{2}\Smodel\).

The passage from this local algebra to a global base is not automatic.
We exhibit a theorem-level non-derivability result: identical local pieces
can be assembled by different incidence data into non-homeomorphic quotient
spaces.  We therefore isolate explicit base-gluing data and prove that,
under a closed-gluing-graph condition, countability, and a sufficient
smoothness condition, the quotient is a Hausdorff second-countable smooth
four-manifold \(M\).

A second non-derivability result separates the internal Lorentz transition
system from the genuine tangent transition system \(D\phi_{ij}\).
Their identification requires an additional regular solder datum.  Conditional
on such a solder, the intrinsic form on \(\Smodel\) descends to a smooth,
chart-independent, nondegenerate Lorentzian \((0,2)\)-tensor of signature
\((1,3)\) on \(TM\); a project-native orientation reduction and a coherent
future-cone reduction follow.  The same solder identifies the tangent
Lorentz cocycle with the projection of a native Spin cocycle, forcing the
project-native orientation and Spin-lifting obstructions to vanish.
No identification of these project-native obstructions with
\(w_{1}(TM)\) or \(w_{2}(TM)\) is made.

Finally, on one fixed periodic control, we analyse a selected one-parameter
Spin transition family generated by a chosen Clifford direction.  Every
finite parameter value is regularly admissible; all parameters lie in one
full Spin \Cech\ gauge orbit; and the complete regular-solder solution spaces
at two parameter values are explicitly bijective with identical induced
tangent metrics under the correspondence.  This is a scoped negative control
on a transition-cocycle deformation, not a general deformation
classification and not a construction of connection, holonomy, curvature,
or dynamics.
\end{abstract}
\clearpage

\tableofcontents
\clearpage

% ======================================================================
\section{Introduction}
% ======================================================================

The construction of a Lorentzian spin geometry is usually formulated after
a smooth manifold, a Lorentz metric, and the associated frame bundle have
already been supplied; the Spin structure then appears as a lift of an
orthonormal frame bundle subject to a topological obstruction
\cite{LawsonMichelsohn,MilnorStasheff,Steenrod}.  Here the order of
dependence is deliberately reversed.  The starting object is only
\[
   V=\R^{3},\qquad
   \langle u,v\rangle=u_{1}v_{1}+u_{2}v_{2}+u_{3}v_{3},
\]
with no four-dimensional base space, tangent bundle, Lorentz form,
orientation, Spin bundle, connection, or curvature in the input.

There are then three logically separate reconstruction questions.

First, what local Lorentzian and intrinsic Spin-algebraic structures can be
constructed from the local Euclidean algebra?  Second, what information must be
added before local copies of the resulting model can be assembled into a
global smooth manifold?  Third, after such a manifold has been obtained,
what additional datum is required to identify the internal Lorentz system
with its genuine tangent geometry?

The separation of these questions is essential.  A local model is not a
manifold; fibre transition data do not determine which local base points
are to be identified; and an internal Lorentz cocycle is not, without
further structure, the derivative cocycle of a manifold atlas.  We make
each of these boundaries explicit and place the positive reconstruction
theorems immediately after the relevant missing datum.

The local algebra belongs to the standard circle of ideas around real
Clifford algebras, Spin groups, spin factors and Lorentz cones
\cite{Chevalley,LawsonMichelsohn,JordanNeumannWigner,FarautKoranyi}.
The global steps use standard quotient, manifold, bundle and \Cech\ ideas
\cite{LeeSmooth,Husemoller,Steenrod,KobayashiNomizu}, but the point of the
present organization is not to assume the usual global Lorentzian package
in advance.  Instead, the precise dependency chain is
\[
\boxed{
\begin{array}{c}
(\R^{3},\langle\cdot,\cdot\rangle)
\\[0.25em]\Downarrow\\[-0.15em]
(\Smodel,N,B,\Cl_{3},\SpinG,\rho,\gB)
\\[0.25em]
+\;\text{base incidence and gluing}
\\[-0.05em]\Downarrow\\[-0.15em]
\text{topological four-manifold candidate}
\\[0.25em]
+\;\text{closed gluing graph + countable index + smooth gluing}
\\[-0.05em]\Downarrow\\[-0.15em]
M^{4},\quad TM,\quad D\phi_{ij}
\\[0.25em]
+\;\text{regular solder}
\\[-0.05em]\Downarrow\\[-0.15em]
(TM,g,\text{orientation reduction},
\text{ future-cone reduction},
\text{ native Spin reconvergence}).
\end{array}}
\tag{1.1}
\]
The boxes marked by plus signs are genuine gates: their inhabitance is not
derived from the preceding local algebra.

The following six statements organize the paper.

\begin{theorem}[Local reconstruction]
\label{thm:intro-local}
From the specified Euclidean datum \(V=\R^{3}\), the construction produces:
\begin{enumerate}
\item a four-dimensional real spin factor
      \(\Smodel=\R\oplus V\) with Lorentz form of signature \((1,3)\);
\item the real Clifford algebra \(\Cl_{3}=\Cl(V,q)\), with the paravector
      embedding of \(\Smodel\);
\item an intrinsic Spin group \(\SpinG\), an intrinsic proper
      cone-preserving Lorentz group \(\LorG\), and a continuous surjective
      homomorphism \(\rho:\SpinG\to\LorG\) whose central kernel is exactly
      \(\{\pm1\}\), together with continuous local sections;
\item a six-dimensional Lorentz-skew Lie algebra \(\gB\) and a linear
      equivalence \(\Lambda^{2}\Smodel\simeq\gB\).
\end{enumerate}
\end{theorem}

\begin{theorem}[Conditional smooth globalization]
\label{thm:intro-global}
The local model does not determine a global base.  Given explicit compatible
base-incidence/gluing data \(B\), a closed gluing graph, a countable index
set, and the sufficient condition that the base transition maps are
\(C^\infty\), the quotient \(M_B\) is a Hausdorff second-countable smooth
four-manifold modelled on \(\Smodel\).
\end{theorem}

\begin{theorem}[Tangent--internal independence]
\label{thm:intro-independence}
For the reconstructed smooth manifold, the tangent transition system is the
derivative system \(D\phi_{ij}\).  A coherent native Spin cocycle gives,
independently, an internal Lorentz transition system
\(\rho(\widetilde g_{ij})\).  Smooth base gluing together with coherent
native Spin data does not force these two systems to agree.
\end{theorem}

\begin{theorem}[Regular solder closure]
\label{thm:intro-solder}
A regular solder is an additional smooth local family of fibrewise linear
equivalences satisfying the exact overlap intertwining law between
\(D\phi_{ij}\) and \(\rho(\widetilde g_{ij})\).
Conditional on its existence, it yields a bundle-level equivalence
certificate between the internal Lorentz bundle and \(TM\), with continuous
total-space maps in both directions and smooth local representatives.  It
also induces a smooth Lorentzian tangent metric of signature \((1,3)\).
\end{theorem}

\begin{theorem}[Global reconvergence]
\label{thm:intro-reconvergence}
For every regular solder, the tangent Lorentz transition system in solder
frames equals the projected native Spin cocycle.  Consequently the
project-native determinant obstruction is a coboundary and the
project-native fixed-cover Spin-lifting obstruction is trivial.  A regular
solder also yields a coherent future-cone reduction.
\end{theorem}

\begin{theorem}[Selected-control gauge absorption]
\label{thm:intro-deformation}
On the fixed periodic control and for the selected one-parameter native Spin
transition family considered in \cref{sec:deformation}, every finite
parameter value admits a regular solder.  Any two parameter values define
full-Spin-\Cech-gauge-equivalent cocycles.  Their complete regular-solder
solution spaces are explicitly bijective, and corresponding solders induce
the same tangent Lorentz metric.
\end{theorem}

The hypotheses in \cref{thm:intro-global,thm:intro-solder} are not cosmetic.
They are exactly where the construction ceases to be intrinsic to the
initial Euclidean local data.  In particular, this paper does \emph{not}
prove an implication
\[
   \R^{3}\quad\Longrightarrow\quad\text{a canonical spacetime}.
\]
It proves the local implication in \cref{thm:intro-local}, identifies the
additional global data required to continue, and then proves what follows
from those data.

% ======================================================================
\clearpage
\section{The Euclidean primitive and the Lorentzian spin factor}
\label{sec:local}
% ======================================================================

\subsection{Primitive input}

Throughout the local part of the paper,
\[
   V:=\R^{3}
\]
carries only its standard positive-definite bilinear form
\[
   \langle u,v\rangle
   :=u_{1}v_{1}+u_{2}v_{2}+u_{3}v_{3}.
\]
Set
\[
   q_{3}(v):=\langle v,v\rangle.
\]
No Lorentzian structure has yet been introduced.

\begin{definition}[Real spin factor]
\label{def:spin-factor}
Define
\[
   \Smodel:=\R\oplus V
\]
and, for \(x=(a,u)\), \(y=(b,v)\), define
\[
   x\circ y
   :=
   \bigl(ab+\langle u,v\rangle,\; av+bu\bigr).
\]
The distinguished unit is
\[
   \mathbf{1}_{\Smodel}:=(1,0).
\]
\end{definition}

The product is bilinear and commutative.  As an auxiliary local identity
for this chosen spin-factor construction one also has the Jordan identity
\[
   (x\circ y)\circ(x\circ x)
   =
   x\circ\bigl(y\circ(x\circ x)\bigr).
\]
Thus \((\Smodel,\circ,\mathbf{1}_{\Smodel})\) is the rank-two real spin
factor; compare the standard Jordan-algebra treatment in
\cite{JordanNeumannWigner,FarautKoranyi}.  This local identity is supporting
algebra, not an additional global reconstruction claim.

\begin{definition}[Intrinsic quadratic and bilinear forms]
\label{def:NSBS}
For \(x=(a,u)\in\Smodel\), set
\[
   N(x):=a^{2}-\langle u,u\rangle.
\]
Its polarization is
\[
   B(x,y)
   :=\frac{N(x+y)-N(x)-N(y)}{2}
   =ab-\langle u,v\rangle
\]
for \(y=(b,v)\).
\end{definition}

\begin{proposition}[Lorentz signature]
\label{prop:signature}
The form \(B\) is symmetric and nondegenerate, and has signature \((1,3)\).
In the standard basis
\[
   \mathbf{1}_{\Smodel},\quad
   (0,\varepsilon_{0}),\quad(0,\varepsilon_{1}),\quad(0,\varepsilon_{2}),
\]
its Gram matrix is
\[
   \operatorname{diag}(1,-1,-1,-1).
\]
\end{proposition}

\begin{proof}
The displayed formula
\(B((a,u),(b,v))=ab-\langle u,v\rangle\)
gives the Gram matrix directly.  Since the Euclidean form on \(V\) is
positive definite, the first basis direction is positive and the remaining
three are negative.  The determinant is nonzero, hence \(B\) is
nondegenerate.
\end{proof}

A second auxiliary local identity of the spin-factor construction is the
Cayley--Hamilton identity
\[
   x\circ x-\operatorname{tr}_{\Smodel}(x)\,x
   +N(x)\mathbf{1}_{\Smodel}=0,
   \qquad
   \operatorname{tr}_{\Smodel}(a,u)=2a.
\tag{2.1}
\]
The square cone
\[
   C_{\Smodel}:=
   \{x\in\Smodel:\exists y\in\Smodel,\ x=y\circ y\}
\tag{2.2}
\]
provides an intrinsic time-cone structure at the level of the local model.

\begin{definition}[Intrinsic Lorentz group]
\label{def:GLor}
Let \(\LorG\) be the group of real-linear self-equivalences
\(F:\Smodel\to\Smodel\) satisfying
\[
   N(Fx)=N(x),\qquad
   x\in C_{\Smodel}\Longleftrightarrow Fx\in C_{\Smodel},
   \qquad
   \det F=1.
\tag{2.3}
\]
\end{definition}

The definition is internal to \((\Smodel,N,C_{\Smodel})\).  In particular,
no matrix realization over \(\mathbb C\), no pre-existing Lorentz group on
a spacetime, and no frame bundle enters the construction.

\begin{boundary}[Local means local]
\label{bdry:no-manifold-yet}
At this point \(\Smodel\) is a four-dimensional real vector space with an
intrinsic Lorentz form.  It is not a tangent space and there is no manifold,
chart, atlas, base point, or covering space in the construction.
\end{boundary}

% ======================================================================
\clearpage
\section{Clifford envelope and the intrinsic Spin double cover}
\label{sec:clifford}
% ======================================================================

\subsection{The Euclidean Clifford algebra}

Let
\[
   \Cl_{3}:=\Cl(V,q_{3})
\]
be the real Clifford algebra characterized by the relation
\[
   \iota(v)^{2}=q_{3}(v)\mathbf{1}.
\tag{3.1}
\]
For the standard Euclidean basis
\(\varepsilon_{0},\varepsilon_{1},\varepsilon_{2}\) of \(V\), write
\(\gamma_{i}:=\iota(\varepsilon_{i})\).  Then
\[
   \gamma_{i}^{2}=1,\qquad
   \gamma_{i}\gamma_{j}=-\gamma_{j}\gamma_{i}
   \quad(i\neq j).
\tag{3.2}
\]
The signs are those of the \emph{positive} Euclidean quadratic form on
\(V\); the Lorentzian signature has already appeared in the distinct
four-dimensional spin factor \(\Smodel\).

\begin{definition}[Paravector embedding]
\label{def:paravector}
Define
\[
   A:\Smodel\longrightarrow\Cl_{3},
   \qquad
   A(a,v):=a\,\mathbf{1}+\iota(v).
\tag{3.3}
\]
\end{definition}

Let \(c\) denote the Clifford conjugation used in the construction.  Its
restriction to paravectors changes the sign of the vector part.  The key
identity is
\[
   A(x)c(A(x))=N(x)\mathbf{1}.
\tag{3.4}
\]

\begin{proposition}[Jordan--Clifford compatibility]
\label{prop:jordan-clifford}
For all \(x,y\in\Smodel\),
\[
   A(x\circ y)
   =
   \frac12\bigl(A(x)A(y)+A(y)A(x)\bigr).
\tag{3.5}\label{eq:jordan-clifford}
\]
Moreover, the image \(A(\Smodel)\) generates \(\Cl_{3}\) as a real
algebra.
\end{proposition}

\begin{proof}
Expand the product using
\[
   \iota(u)\iota(v)+\iota(v)\iota(u)
   =2\langle u,v\rangle\mathbf{1}.
\]
The scalar part becomes
\(2(ab+\langle u,v\rangle)\), while the vector part becomes
\(2(av+bu)\), proving \eqref{eq:jordan-clifford}.  Since \(A(0,v)=\iota(v)\), the image
contains every Clifford generator; the universal generation property of
the Clifford algebra then gives the second assertion.
\end{proof}

As standard background for the Mathlib \texttt{CliffordAlgebra}
construction used here, \(\Cl_{3}\) carries the usual universal property:
any real algebra receiving a linear map \(f:V\to\mathcal A\) with
\(f(v)^{2}=q_{3}(v)\mathbf1\) receives a unique algebra homomorphism from
\(\Cl_{3}\).  This is a property of the standard Clifford-algebra
construction rather than an additional registered global claim of the
present reconstruction; see \cite{Chevalley,LawsonMichelsohn}.

\subsection{Intrinsic Spin elements}

\begin{definition}[Spin element and Spin group]
\label{def:spingroup}
An element \(g\in\Cl_{3}\) is a Spin element when
\[
    g\,c(g)=c(g)\,g=1.
\tag{3.6}
\]
The corresponding units form a subgroup
\[
    \SpinG\subset\Cl_{3}^{\times}.
\tag{3.7}
\]
\end{definition}

The Clifford action used in the construction preserves the paravector
subspace and defines a real-linear action
\[
   \operatorname{spinLor}(g):\Smodel\longrightarrow\Smodel.
\]
For every \(g\in\SpinG\), this action belongs to \(\LorG\).
Hence
\[
   \rho:\SpinG\longrightarrow\LorG,
   \qquad
   \rho(g)=\operatorname{spinLor}(g),
\tag{3.8}
\]
is a group homomorphism.

\begin{theorem}[Intrinsic double cover]
\label{thm:double-cover}
The homomorphism \(\rho\) is surjective and
\[
   \ker\rho=\{1,-1\}.
\tag{3.9}
\]
The kernel is central and has two elements.  With the topologies inherited
from the finite-dimensional Clifford model and the induced topology on
\(\LorG\), both groups are Hausdorff topological groups, \(\rho\) is
continuous, the kernel is discrete, and every \(k\in\LorG\) has an open
neighbourhood \(U\) admitting a continuous local section
\[
   s:U\longrightarrow\SpinG,
   \qquad
   \rho\circ s=\id_{U}.
\tag{3.10}
\]
\end{theorem}

\begin{proof}
We separate surjectivity, the kernel calculation, and the local-section
statement.

\emph{Surjectivity.}  Let \(F\in\LorG\) and put \(p=F(1_{\Smodel})\).
Because \(F\) preserves both \(N\) and the intrinsic square cone,
\[
    N(p)=1,
    \qquad
    p\in C_{\Smodel},
    \qquad
    p_0\ge 0.
\]
Write \(p=(a,u)\).  From \(a^2-\langle u,u\rangle=1\) and \(a\ge0\) one
has \(a\ge1\).  Hence the forward spin-factor square root
\[
 r=\sqrt[\Smodel]{p}
   :=\left(\sqrt{\frac{a+1}{2}},
      \frac{u}{2\sqrt{(a+1)/2}}\right)
\]
is well defined and satisfies \(r\circ r=p\) and \(N(r)=1\).  Set
\(\beta=A(r)\in\Cl_3\).  The Clifford norm identity makes \(\beta\) a Spin
element, and the Jordan--Clifford compatibility gives
\[
   \rho(\beta)(1_{\Smodel})=p.
\]
Consequently
\[
   G:=\rho(\beta)^{-1}F
\]
fixes \(1_{\Smodel}\).  The stabilizer of \(1_{\Smodel}\) acts on the
spatial summand as an element \(R\in SO(3)\).  By Cartan--Dieudonn\'e,
\(R\) is a product of hyperplane reflections in unit spatial vectors.
For a unit vector \(u\), the Clifford generator \(\iota(u)\) acts on the
spatial summand as minus the corresponding hyperplane reflection.  Since
\(\det R=1\), an even number of reflections occurs; the accumulated minus
sign therefore disappears, and the product of the corresponding Clifford
generators is a Spin element \(\gamma\) with \(\rho(\gamma)=G\).  Thus
\(\rho(\beta\gamma)=F\).

\emph{Kernel.}  Suppose \(g\in\SpinG\) acts trivially on every paravector.
Evaluating at \(1_{\Smodel}\) first shows that Clifford reversion and
Clifford conjugation agree on \(g\).  The trivial-action equation then
implies
\[
    gA(x)=A(x)g
    \qquad\text{for every }x\in\Smodel.
\]
Because the paravector image generates \(\Cl_3\) as a real algebra, \(g\)
commutes with all of \(\Cl_3\).  The intrinsically computed centre is
\(\R\oplus\R\omega\), where
\(\omega=\gamma_0\gamma_1\gamma_2\) and \(\omega^2=-1\).  Hence
\(g=a+b\omega\).  For central elements Clifford conjugation fixes \(g\),
so the Spin normalization gives
\[
   (a+b\omega)^2=1
   \quad\Longrightarrow\quad
   a^2-b^2=1,
   \qquad
   2ab=0.
\]
The first equation excludes \(a=0\); therefore \(b=0\) and
\(a=\pm1\).  Conversely, \(\pm1\) act trivially.  Thus
\(\ker\rho=\{\pm1\}\); centrality and cardinality two follow immediately.

\emph{Continuous local sections.}  The construction is also intrinsic.  Define the Clifford contraction
\[
   \mathcal F(y)
   :=y+\sum_{i=0}^{2}\gamma_i y\gamma_i .
\]
In the eight-monomial Clifford frame this kills the six noncentral
monomials and multiplies the two central ones by \(4\), so
\(\mathcal F(y)\in\R\oplus\R\omega\).  Next define the reconstruction map
\[
   \mathcal R(F)
   :=A(F1_{\Smodel})
     +\sum_{i=0}^{2}A\!\left(F(0,\varepsilon_i)\right)\gamma_i .
\]
For a Spin element \(g\), direct Clifford multiplication gives
\[
   \mathcal R(\rho(g))
   =g\,\mathcal F(\operatorname{rev}g),
\]
so reconstruction recovers \(g\) up to a central factor.  Identify the
centre \(\R\oplus\R\omega\) with \(\mathbb C\) by the real-algebra embedding
\(\operatorname{cx}:\mathbb C\to\Cl_3\) determined by
\(\operatorname{cx}(i)=\omega\), and choose the linear retraction
\(\operatorname{cxRetract}\) used by the reconstruction.  Define the
continuous discriminant
\[
   \Delta(F)
   :=\operatorname{cxRetract}\!\left(
      \mathcal R(F)\,c(\mathcal R(F))
   \right).
\]
The key identity is
\[
   \mathcal R(\rho(g))=\operatorname{cx}(w)g
   \quad\Longrightarrow\quad
   \Delta(\rho(g))=w^2.
\]
Thus the central ambiguity is visible only through a square.  At the identity,
\(\Delta(1)=16\).  On the open neighbourhood where \(\Delta(F)\) lies in
a slit plane avoiding the branch cut, the principal complex square root is
continuous.  Dividing the reconstructed Clifford element by this square root
leaves exactly the two choices differing by \(\{\pm1\}\), and yields a
continuous Spin-valued section \(s\) with \(\rho(s(F))=F\) and \(s(1)=1\).
Finally, for arbitrary \(k\in\LorG\), choose one lift \(u_k\) of \(k\) and
left-translate the identity neighbourhood and its section.  This produces
an open neighbourhood of \(k\) with a continuous local section.  A
different choice of \(u_k\) changes that section only by the central
kernel element \(-1\).
\end{proof}

Thus one obtains the exact central extension
\[
   1\longrightarrow\{\pm1\}
   \longrightarrow\SpinG
   \overset{\rho}{\longrightarrow}\LorG
   \longrightarrow1.
\tag{3.11}
\]
A classical complex-matrix description may be used as a comparison
endpoint, but it is not an input to \cref{thm:double-cover}.  In particular,
the construction does not define \(\SpinG\) by first declaring it to be
\(\mathrm{SL}(2,\mathbb C)\).

% ======================================================================
\clearpage
\section{The intrinsic infinitesimal Lorentz sector}
\label{sec:lie}
% ======================================================================

The local group construction has a parallel infinitesimal object that can
be defined without first making \(\SpinG\) into a Lie group.

\begin{definition}[Lorentz-skew endomorphisms]
\label{def:gB}
Let
\[
   \gB
   :=
   \bigl\{
      T\in\End_{\R}(\Smodel):
      B(Tx,y)+B(x,Ty)=0
      \ \text{for all }x,y\in\Smodel
   \bigr\}.
\tag{4.1}
\]
\end{definition}

For \(u,v\in\Smodel\), define
\[
   K_{u,v}(x)
   :=
   B(v,x)u-B(u,x)v.
\tag{4.2}\label{eq:Kuv}
\]
Then \(K_{u,v}\in\gB\), the map is alternating in \(u,v\), and it therefore
induces a linear map
\[
   \Phi:\Lambda^{2}\Smodel\longrightarrow\gB,
   \qquad
   u\wedge v\longmapsto K_{u,v}.
\tag{4.3}
\]

\begin{theorem}[Bivector realization]
\label{thm:bivector}
The map \(\Phi\) is a linear equivalence.  Consequently
\[
   \dim_{\R}\gB
   =
   \dim_{\R}\Lambda^{2}\Smodel
   =6.
\tag{4.4}
\]
Moreover, \(\gB\) is closed under the commutator
\[
   [A,C]:=AC-CA,
\tag{4.5}
\]
so it is a six-dimensional real Lie algebra.
\end{theorem}

\begin{proof}
Skewness of \(K_{u,v}\) follows by expanding \eqref{eq:Kuv} and using the
symmetry of \(B\).  Closure of \(\gB\) under the commutator follows by
applying the skewness identity twice:
\[
  B(ACx,y)=-B(Cx,Ay)=B(x,CAy),
\]
and similarly with \(A\) and \(C\) interchanged.

It remains to prove that \(\Phi\) is bijective.  Let
\[
   \tau=(1,0),
   \qquad
   b_i=(0,\varepsilon_i),\quad i=0,1,2,
\]
where \(\varepsilon_0,\varepsilon_1,\varepsilon_2\) is the fixed Euclidean
basis of \(V\).  For \(A\in\gB\), introduce the six coefficients
\[
\begin{aligned}
   c_0&=(A\tau)_V^{0},&
   c_1&=(A\tau)_V^{1},&
   c_2&=(A\tau)_V^{2},\\
   c_3&=\langle (Ab_0)_V,\varepsilon_1\rangle,&
   c_4&=\langle (Ab_0)_V,\varepsilon_2\rangle,&
   c_5&=\langle (Ab_1)_V,\varepsilon_2\rangle.
\end{aligned}
\tag{4.6a}\label{eq:gb-six-coordinates}
\]
The \(B\)-skew condition determines every other matrix coefficient from
these six: it forces \((A\tau)_0=0\), determines the scalar component of
\(A(0,u)\) by \(\langle(A\tau)_V,u\rangle\), and makes the spatial block
antisymmetric.  Hence if all six numbers in
\eqref{eq:gb-six-coordinates} vanish, then \(A\tau=0\) and
\(Ab_i=0\) for \(i=0,1,2\).  Since \(\tau,b_0,b_1,b_2\) span
\(\Smodel\), one has \(A=0\).  Thus the six coordinates are rigid.

Conversely, for \(c=(c_0,\ldots,c_5)\in\R^6\), define the explicit
bivector
\[
\begin{aligned}
   z(c):={}&-c_0\,\tau\wedge b_0
            -c_1\,\tau\wedge b_1
            -c_2\,\tau\wedge b_2\\
          &+c_3\,b_0\wedge b_1
            +c_4\,b_0\wedge b_2
            +c_5\,b_1\wedge b_2 .
\end{aligned}
\tag{4.6b}\label{eq:zmap}
\]
A direct evaluation of \(K_{u,v}\) on the four basis vectors shows that
the six coordinates of \(\Phi(z(c))\) are exactly \(c_0,\ldots,c_5\).
For any \(A\in\gB\), taking \(c\) to be the six coordinates of \(A\)
therefore gives an endomorphism \(\Phi(z(c))\) with the same six
coordinates as \(A\); rigidity yields \(\Phi(z(c))=A\).  Thus \(\Phi\) is
surjective.

For injectivity, the map \(z:\R^6\to\Lambda^2\Smodel\) in
\eqref{eq:zmap} is injective because the six coordinates of
\(\Phi(z(c))\) recover \(c\).  Since
\[
   \dim\R^6=6=\dim\Lambda^2\Smodel,
\]
this injective linear map is also surjective.  If \(\Phi(z)=\Phi(z')\),
write \(z=z(c)\) and \(z'=z(c')\); applying the six coordinate functionals
gives \(c=c'\), hence \(z=z'\).  Therefore \(\Phi\) is injective as well,
and so it is a linear equivalence.  The dimension statement follows.
\end{proof}

\begin{boundary}[The first open infinitesimal edge]
\label{bdry:lie-spin}
The statement
\[
   \operatorname{Lie}(\SpinG)\simeq\gB
\tag{4.7}
\]
is not part of the present construction.  The intrinsic \(\SpinG\) is
currently certified as a topological group, not as a smooth Lie group.
Accordingly, the differential
\[
   d\rho_{e}:\operatorname{Lie}(\SpinG)
   \longrightarrow\gB
\tag{4.8}
\]
has not been constructed.  These are theorem-level gaps between the
group-level double cover and the already available Lie algebra \(\gB\).
\end{boundary}

% ======================================================================
\clearpage
\section{From local model to global base}
\label{sec:globalization}
% ======================================================================

The next step is deliberately not phrased as a consequence of the Spin
group.  The reason is structural: a fibre transition tells us how to
compare fibres over a point already known to lie in two patches.  It does
not say which points of two local copies represent the same point of a
global base.

\subsection{Base-gluing data}

Let \(I\) be an index set.  A local piece is an open subset
\(D_i\subset\Smodel\).

\begin{definition}[Base-gluing datum]
\label{def:basegluing}
A base-gluing datum \(B\) consists of:
\begin{enumerate}
\item open sets \(D_i\subset\Smodel\);
\item open incidence domains \(W_{ij}\subset D_i\);
\item continuous maps
      \(\phi_{ij}:\Smodel\to\Smodel\), used on \(W_{ij}\);
\end{enumerate}
such that, on the indicated domains,
\[
\begin{aligned}
   &W_{ii}=D_i,
   &&\phi_{ii}(x)=x,\\
   &\phi_{ij}(W_{ij})\subset W_{ji},
   &&\phi_{ji}(\phi_{ij}(x))=x,
\end{aligned}
\tag{5.1}
\]
and whenever \(x\in W_{ij}\) and
\(\phi_{ij}(x)\in W_{jk}\),
\[
   x\in W_{ik},
   \qquad
   \phi_{jk}(\phi_{ij}(x))=\phi_{ik}(x).
\tag{5.2}
\]
\end{definition}

This is an explicit \emph{sufficient} base-gluing primitive.  We do not
claim that every field in Definition~\ref{def:basegluing} is separately necessary,
or that this is the unique minimal encoding of incidence information.

Form the tagged disjoint union
\[
   X_{0}:=\coprod_{i\in I}D_i
\tag{5.3}
\]
and define
\[
   (i,x)\sim(j,y)
   \quad\Longleftrightarrow\quad
   x\in W_{ij}\ \text{ and }\ \phi_{ij}(x)=y.
\tag{5.4}
\]
The axioms above make \(\sim\) an equivalence relation.  Define
\[
   M_B:=X_{0}/\!\sim
\tag{5.5}
\]
with the quotient topology.

\begin{proposition}[Topological reconstruction]
\label{prop:top-reconstruct}
The quotient map \(X_{0}\to M_B\) is open.  Each local piece \(D_i\)
embeds homeomorphically onto an open subset of \(M_B\), and these images
cover \(M_B\).  Thus \(M_B\) carries a charted-space structure modelled on
\(\Smodel\).
\end{proposition}

\begin{proof}
The saturation of an open subset of \(X_{0}\) is open because each
incidence set \(W_{ij}\) is open and \(\phi_{ij}\) is a homeomorphism
between \(W_{ij}\) and \(W_{ji}\), with inverse \(\phi_{ji}\).  Hence the
quotient map is open.  Equality of two points from the same tagged piece
forces equality of their local coordinates by \(\phi_{ii}=\id\), so the
map from each \(D_i\) is injective.  Openness and continuity then give an
open embedding.  Surjectivity of the quotient map proves that the chart
images cover.
\end{proof}

Hausdorffness is not a formal consequence of Definition~\ref{def:basegluing}; the
usual line-with-two-origins phenomenon already shows why a separation
condition is needed.

\begin{definition}[Closed gluing graph]
\label{def:closedgraph}
The gluing graph of \(B\) is closed when
\[
   \Gamma_B
   :=
   \{(p,q)\in X_0\times X_0:p\sim q\}
\tag{5.6}
\]
is closed.
\end{definition}

\begin{proposition}[Hausdorffness and second countability]
\label{prop:haus-second}
If \(\Gamma_B\) is closed, then \(M_B\) is Hausdorff.  If, in addition,
\(I\) is countable, then \(M_B\) is second countable.
\end{proposition}

\begin{proof}
Let \(q:X_0\to M_B\) be the quotient map.  By
Proposition~\ref{prop:top-reconstruct}, \(q\) is continuous, open, and surjective.
Consequently
\[
   q\times q:X_0\times X_0\longrightarrow M_B\times M_B
\]
is again an open quotient map.  The preimage of the diagonal is exactly the
gluing relation:
\[
   (q\times q)^{-1}(\Delta_{M_B})
   =\{(p,q):p\sim q\}
   =\Gamma_B.
\]
If \(\Gamma_B\) is closed, the quotient property therefore implies that
\(\Delta_{M_B}\) is closed.  This is equivalent to Hausdorffness.

For second countability, each \(D_i\) is an open subspace of the
finite-dimensional real vector space \(\Smodel\), hence is second
countable.  If \(I\) is countable, the tagged disjoint union
\(X_0=\coprod_i D_i\) is second countable.  An open quotient of a
second-countable space is second countable: the images under \(q\) of a
countable basis form a countable basis of \(M_B\).  Hence \(M_B\) is second
countable.
\end{proof}

\subsection{Why the global base is additional}

The need for incidence data is itself a theorem, not merely a warning.

\begin{theorem}[Local pieces do not determine the base]
\label{thm:base-nonderivability}
There exist two base-gluing data over the same two indexed local domains
such that the local pieces agree exactly but their quotient bases are not
homeomorphic.
\end{theorem}

\begin{proof}
Take a nonempty connected open local domain \(D\subset\Smodel\) and two
copies indexed by \(\{\mathsf{false},\mathsf{true}\}\).  In the first
gluing, make every pair of pieces fully incident and use the identity
identification.  In the second, take \(W_{ij}=\varnothing\) for \(i\neq j\)
and again use the identity where defined.  The first quotient is connected;
the second is the disjoint union of two nonempty open components and is not
connected.  Hence no homeomorphism can exist between them.
\end{proof}

\begin{corollary}
\label{cor:fibre-not-base}
No datum that only prescribes a fibre transition law over already specified
overlaps can, by itself, determine the missing cross-piece incidence of the
base.
\end{corollary}

\subsection{The smooth gate}

For the reconstructed atlas, the coordinate change from piece \(i\) to
piece \(j\) is precisely \(\phi_{ij}\) on \(W_{ij}\).

\begin{definition}[Smooth gluing]
\label{def:smoothgluing}
A base-gluing datum is smoothly compatible when
\[
   \phi_{ij}|_{W_{ij}}
\]
is \(C^\infty\) for every \(i,j\).
\end{definition}

\begin{theorem}[Conditional smooth four-manifold]
\label{thm:smooth-manifold}
Suppose:
\begin{enumerate}
\item \(B\) is a base-gluing datum over \(\Smodel\);
\item its gluing graph is closed;
\item the index set \(I\) is countable;
\item \(B\) satisfies Definition~\ref{def:smoothgluing}.
\end{enumerate}
Then \(M_B\) is Hausdorff, second countable, and a \(C^\infty\)
four-manifold modelled on \(\Smodel\).
\end{theorem}

\begin{proof}
The Hausdorff and second-countable assertions are exactly
Proposition~\ref{prop:haus-second}.  It remains to identify the smooth structure of
the actual quotient atlas.

For each inhabited piece \(D_i\), the open embedding of
Proposition~\ref{prop:top-reconstruct} yields a chart whose target is \(D_i\).  On an
overlap of the \(i\)- and \(j\)-charts, the source of the coordinate change
is not merely contained in the primitive incidence set: it is exactly
\(W_{ij}\).  If \(y\in W_{ij}\), the quotient relation identifies
\((i,y)\) with \((j,\phi_{ij}(y))\), and therefore the coordinate change is
literally
\[
     y\longmapsto\phi_{ij}(y).
\]
Thus every transition map of the reconstructed atlas is a restriction of
one of the functions already required to be \(C^\infty\) by
Definition~\ref{def:smoothgluing}.  The atlas consequently takes values in the
\(C^\infty\) structure groupoid, so \(M_B\) is a smooth manifold.
Finally, the model space is
\(\Smodel=\R\oplus\R^3\), hence \(\dim_{\R}\Smodel=4\).

Notice that the argument proves sufficiency of the smooth-gluing
hypothesis for this reconstructed atlas.  It does not claim the converse or
necessity of every primitive smoothness field.
\end{proof}

\begin{control}[Smoothness is not automatic]
\label{control:nonsmooth}
There are base-gluing data satisfying the topological gluing axioms whose
transition is a homeomorphism but not differentiable at the origin.  An
explicit control is
\[
   H(t,x)_0=t,
   \qquad
   H(t,x)_k=x_k+|t| \quad(k=1,2,3),
\]
with inverse
\[
   H^{-1}(t,x)_0=t,
   \qquad
   H^{-1}(t,x)_k=x_k-|t|.
\]
Both maps are continuous, so the topological reconstruction applies.  But
along the line \(t\mapsto(t,0)\), every shifted spatial component contains
\(|t|\), which is not differentiable at \(t=0\).  Hence the smooth-gluing
condition fails.
\end{control}

\begin{remark}[Sufficiency, not necessity]
\label{rem:smooth-sufficient}
For the specific quotient-atlas packaging used here, smooth gluing is
proved sufficient.  A converse is not asserted.  In particular, quotient
representative selection may cause the reconstructed atlas to omit some
index-labelled chart even when the corresponding primitive
\(\phi_{ij}\) is non-smooth.
\end{remark}

% ======================================================================
\clearpage
\section{Internal Spin data and the genuine tangent cocycle}
\label{sec:tangent}
% ======================================================================

Fix from now on a smooth base gluing \(B\) satisfying the hypotheses needed
to obtain \(M=M_B\).

\subsection{Genuine tangent transitions}

Because the actual atlas transition is \(\phi_{ij}\), the tangent bundle
transition is its derivative.

\begin{proposition}[Tangent transition law]
\label{prop:tangent-transition}
For \(y\in W_{ij}\), the genuine tangent coordinate change of the
reconstructed manifold is
\[
   T_{ij}(y):=D\phi_{ij}(y):
   \Smodel\longrightarrow\Smodel.
\tag{6.1}\label{eq:tangent-transition}
\]
It is a linear isomorphism, with inverse
\[
   D\phi_{ji}(\phi_{ij}(y)).
\tag{6.2}\label{eq:tangent-inverse}
\]
The transitions satisfy the derivative cocycle law by the chain rule.
\end{proposition}

The invertibility in \eqref{eq:tangent-inverse} is derived from
\(\phi_{ji}\circ\phi_{ij}=\id\) on the open incidence domain rather than
inserted as a separate hypothesis.

\begin{remark}[A type-level shortcut is not a solder]
In the formal implementation, tangent spaces of a modelled manifold are
represented using the model vector space, so a type-level identification
with \(\Smodel\) is available.  This does not identify transition systems
and plays no role in the geometric coupling below.  The invariant datum is
\eqref{eq:tangent-transition}.
\end{remark}

\subsection{Native Spin transitions}

Let \(\cU=\{U_i\}_{i\in I}\) be the cover of \(M\) by the chart images.
A native Spin transition system consists of continuous maps
\[
   \widetilde g_{ij}:U_i\cap U_j\longrightarrow\SpinG
\tag{6.3}
\]
satisfying the \Cech\ cocycle law
\[
   \widetilde g_{ij}(x)\widetilde g_{jk}(x)
   =
   \widetilde g_{ik}(x)
\tag{6.4}
\]
on triple overlaps.  Projection gives the internal Lorentz cocycle
\[
   g_{ij}:=\rho\circ\widetilde g_{ij}:
   U_i\cap U_j\longrightarrow\LorG.
\tag{6.5}
\]
This is a fibre transition system on the same cover, but it is not the
tangent transition system by definition.

\begin{theorem}[Raw transition identification is not forced]
\label{thm:transition-independence}
There exist a smooth two-chart base gluing \(B\) and a coherent native Spin
transition system \(\widetilde g\) such that
\[
   g_{01}(x)=\id_{\Smodel}
   \quad\text{for every }x,
\tag{6.6}
\]
while
\[
   T_{01}(y)=2\,\id_{\Smodel}
   \quad\text{for every }y.
\tag{6.7}
\]
Hence the base and native Spin data do not force
\(T_{ij}=g_{ij}\).
\end{theorem}

\begin{proof}
Take two copies of the whole local model with
\[
   \phi_{01}(y)=2y,\qquad
   \phi_{10}(y)=\tfrac12 y.
\]
All incidence domains are the whole model, and the gluing is smooth.
Thus \(D\phi_{01}=2\,\id\).  Independently, choose the pointwise trivial
native Spin cocycle \(\widetilde g_{ij}=1\); its projection is the identity.
Evaluating at \(\mathbf1_{\Smodel}\) distinguishes the two transformations.
\end{proof}

\begin{remark}
The conclusion is one of non-derivability, not non-existence.  The same
rescaling example admits a compensating regular solder.  Equality of raw
transition matrices is therefore not the correct invariant criterion for
coupling the two bundles.
\end{remark}

% ======================================================================
\clearpage
\section{The regular solder gate}
\label{sec:solder}
% ======================================================================

The missing object is a smoothly varying identification between the
internal Lorentz fibres and tangent fibres that intertwines their
transition laws.

\subsection{Weak fibrewise soldering}

At the weakest level one may choose, for each chart and each base point, a
linear equivalence
\[
   e_i(x):\Smodel\xrightarrow{\;\sim\;}T_xM
\tag{7.1}
\]
such that on \(U_i\cap U_j\),
\[
   e_j(x)v
   =
   e_i(x)\bigl(g_{ij}(x)v\bigr).
\tag{7.2}\label{eq:weak-solder-law}
\]
This algebraic condition is sufficient to transport the intrinsic form
\(B\) to each tangent fibre, but it imposes no regularity on the field
\(x\mapsto e_i(x)\).

\begin{control}[Weak soldering is insufficient]
\label{control:weak}
There exist fibrewise solder data obeying \eqref{eq:weak-solder-law} whose chart
representative is discontinuous.  Consequently fibrewise linear
equivalence alone does not produce a regular vector-bundle comparison or
a smooth tangent metric.
\end{control}

\begin{proof}[Explicit control]
Take the one-chart flat base with identity gluing and the pointwise trivial
Spin cocycle.  In this case \eqref{eq:weak-solder-law} imposes no
nontrivial overlap constraint.  Let \(x_0\) be the chart point represented
by the zero vector and define the frame field
\[
 e(x)=
 \begin{cases}
   2\,\id_{\Smodel}, & x=x_0,\\
   \id_{\Smodel}, & x\ne x_0.
 \end{cases}
\tag{7.2a}\label{eq:bad-solder}
\]
Every fibre map is a linear equivalence, so this is a valid weak solder.
Its local representative, however, jumps from \(\id\) to \(2\id\) at
\(x_0\), and is therefore discontinuous.  The defect is also visible in
the transported quadratic form: evaluated on the intrinsic unit, the
coefficient is \(1\) away from \(x_0\), whereas at \(x_0\) the inverse
frame contributes a factor \(1/2\), giving \(1/4\).  Thus even the induced
fibrewise metric coefficient is discontinuous.  No regular solder can
induce this weak frame field.
\end{proof}

\subsection{Regular local form}

The regular datum is stated in the actual chart coordinates.  Let
\[
   A_i(y)\in\GL(\Smodel)
\]
be a continuous linear automorphism for \(y\in D_i\).

\begin{definition}[Regular solder]
\label{def:regular-solder}
A regular solder for \((B,\widetilde g)\) is a family
\[
   A_i:D_i\longrightarrow\GL(\Smodel)
\tag{7.3}
\]
such that:
\begin{enumerate}
\item the map \(y\mapsto A_i(y)\), viewed in the finite-dimensional space
      of continuous linear maps, is \(C^\infty\) on \(D_i\);
\item on \(y\in W_{ij}\),
\[
  A_j(\phi_{ij}(y))
  =
  T_{ij}(y)\circ A_i(y)\circ
  g_{ij}(x_y),
\tag{7.4}\label{eq:regular-solder-law}
\]
where \(x_y\in M\) is the quotient point represented by \((i,y)\).
\end{enumerate}
\end{definition}

Equation \eqref{eq:regular-solder-law} fixes the orientation of the intertwining square.
The tangent map \(T_{ij}\) converts \(i\)-chart tangent components to
\(j\)-chart tangent components, while the \Cech\ convention
\(e_j=e_i\circ g_{ij}\) converts internal \(j\)-components to internal
\(i\)-components.  The equation is therefore forced by the two pre-existing
conventions rather than selected by analogy.

The inverse family \(A_i(y)^{-1}\) is automatically smooth; no separate
inverse-regularity field is needed.

\begin{theorem}[Regular bundle-equivalence certificate]
\label{thm:bundle-equivalence}
A regular solder determines a fibrewise linear equivalence between the
internal Lorentz bundle defined by \(g_{ij}\) and the genuine tangent bundle
\(TM\), covering \(\id_M\), with:
\begin{enumerate}
\item continuous total-space map and continuous inverse;
\item linear equivalence on each fibre;
\item \(C^\infty\) local representatives \(A_i\);
\item \(C^\infty\) local representatives \(A_i^{-1}\) for the inverse;
\item exact intertwining of the internal and tangent transition systems.
\end{enumerate}
Conversely, a bundle-equivalence certificate with these properties yields a
regular solder.
\end{theorem}

\begin{proof}
Write the internal Lorentz bundle in the local trivializations determined by
the cocycle \(g_{ij}\), and write \(TM\) in the tangent trivializations of
the reconstructed atlas.  On the \(i\)-th chart define the local total-space
map
\[
   (y,v)\longmapsto (y,A_i(y)v).
\tag{7.4a}\label{eq:local-bundle-map}
\]
If \(y\in W_{ij}\), changing the source trivialization by \(g_{ij}\) and
the target trivialization by \(T_{ij}\) gives, by
\eqref{eq:regular-solder-law}, the same point as first applying
\eqref{eq:local-bundle-map} in chart \(i\).  Hence the local maps agree on
equivalent representatives and descend to a globally defined map of total
spaces covering \(\id_M\).  Each fibre map is a linear equivalence because
each \(A_i(y)\) is.

The local representative of the descended map in the \(i\)-th source and
target trivializations is exactly \(A_i\).  Since \(A_i\) is smooth, it is
continuous, and the standard local-trivialization criterion gives
continuity of the total-space map.  Pointwise inversion on
\(\GL(\Smodel)\) is smooth; hence \(y\mapsto A_i(y)^{-1}\) is smooth.  The
same argument applied to these inverse representatives yields a continuous
inverse total-space map.  This proves the five stated properties.

Conversely, suppose such a bundle-equivalence certificate is given.  Read
its local representative in the chart trivializations and call it
\(A_i(y)\).  Fibrewise linear equivalence gives
\(A_i(y)\in\GL(\Smodel)\), the certificate gives \(C^\infty\)-regularity,
and equality of the globally defined map when read in the two overlapping
trivializations is precisely the intertwining identity
\eqref{eq:regular-solder-law}.  Thus the local representatives form a
regular solder.  The two constructions are inverse at the level asserted
in the theorem.
\end{proof}

\begin{remark}[Exact regularity level]
We establish a bundle equivalence with continuous total-space maps in both
directions, fibrewise linear equivalences, and smooth local representatives
for the map and its inverse.  No stronger smooth total-space equivalence is
asserted here.  In particular, the limitation concerns the level at which the
global equivalence is packaged, not the smoothness of the projected Lorentz
transition representatives, which is forced in chart coordinates by a
regular solder.
\end{remark}

Indeed, \eqref{eq:regular-solder-law} can be solved for \(g_{ij}\):
\[
   g_{ij}(x_y)
   =
   A_i(y)^{-1}\circ T_{ij}(y)^{-1}
   \circ A_j(\phi_{ij}(y)),
\tag{7.5}
\]
up to the convention used in \eqref{eq:regular-solder-law}.  Hence regular soldering forces
the projected internal Lorentz transition to possess smooth local
representatives in chart coordinates even though the underlying native
Spin cocycle is only assumed continuous.

\subsection{Gauge freedom of the solder}

For fixed \(B\) and fixed native Spin data, a solder is not canonical.
The regular tangent-gauge group acts on the space of regular solders.

\begin{theorem}[Solder torsor]
\label{thm:solder-torsor}
Whenever the regular-solder solution space for fixed \((B,\widetilde g)\)
is nonempty, the regular tangent-gauge action on it is free and transitive.
\end{theorem}

\begin{proof}
Let \(A=\{A_i\}\) and \(A'=\{A_i'\}\) be two regular solders for the same
base and the same native Spin datum.  Define
\[
   a_i(y):=A_i'(y)\,A_i(y)^{-1}.
\tag{7.6}\label{eq:gauge-between-solders}
\]
The field \(a_i\) is smooth because multiplication and inversion in the
finite-dimensional general linear group are smooth.  Applying the solder
law once to \(A\) and once to \(A'\), and cancelling the common internal
transition, gives
\[
   a_j(\phi_{ij}(y))\,T_{ij}(y)
   =T_{ij}(y)\,a_i(y),
\tag{7.7}\label{eq:tangent-gauge-law}
\]
which is exactly the regular tangent-gauge compatibility law.  By
construction,
\[
   (a\cdot A)_i(y)=a_i(y)A_i(y)=A_i'(y),
\]
so the action is transitive.

For freeness, if \(a\cdot A=a'\cdot A\), then
\(a_i(y)A_i(y)=a_i'(y)A_i(y)\) at every point.  Right cancellation by the
invertible map \(A_i(y)\) gives \(a_i(y)=a_i'(y)\) for all \(i,y\), hence
\(a=a'\).  Pointwise composition and inversion preserve
\eqref{eq:tangent-gauge-law}, so these fields form a group and the action
is free and transitive.
\end{proof}

Thus existence is a genuine gate, while uniqueness is neither expected nor
claimed.

% ======================================================================
\clearpage
\section{The induced tangent Lorentz geometry}
\label{sec:metric}
% ======================================================================

Let \(A=\{A_i\}\) be a regular solder.

\subsection{Metric coefficients}

In chart \(i\), define
\[
   g_i(y)(u,v)
   :=
   B\bigl(A_i(y)^{-1}u,\;A_i(y)^{-1}v\bigr).
\tag{8.1}\label{eq:metric-coeff}
\]

\begin{theorem}[Smooth tangent Lorentz metric]
\label{thm:smooth-metric}
The fields \(g_i\) satisfy:
\begin{enumerate}
\item for fixed \(u,v\in\Smodel\), the function
      \(y\mapsto g_i(y)(u,v)\) is \(C^\infty\) on \(D_i\);
\item \(g_i(y)\) is symmetric and nondegenerate;
\item on every incidence domain,
\[
  g_j(\phi_{ij}(y))
     \bigl(T_{ij}(y)u,T_{ij}(y)v\bigr)
  =
  g_i(y)(u,v);
\tag{8.2}\label{eq:metric-transform}
\]
\item in the \(A_i(y)\)-frame, the form has signature \((1,3)\).
\end{enumerate}
Hence the \(g_i\) satisfy the smooth tensor transformation law and
therefore define a global Lorentzian \((0,2)\)-tensor \(g\) on \(TM\), in
the standard semi-Riemannian sense \cite{ONeill,Nakahara}.  This conclusion
uses the smooth local coefficient fields and their exact tensor
transformation law; no stronger global metric packaging is asserted.
\end{theorem}

\begin{proof}
Smoothness follows from smoothness of \(A_i^{-1}\) and bilinearity of
\(B\).  Symmetry and nondegeneracy are inherited from \(B\) through the
linear equivalence \(A_i(y)\).  For the transformation law, use
\eqref{eq:regular-solder-law} to express the inverse \(j\)-frame after application of
\(T_{ij}\) in terms of the inverse \(i\)-frame and the projected internal
Lorentz transition.  Since \(g_{ij}(x)\in\LorG\) preserves \(B\), these
Lorentz factors cancel and yield \eqref{eq:metric-transform}.  Finally,
\[
   g_i(y)(A_i(y)a,A_i(y)b)=B(a,b),
\]
so Proposition~\ref{prop:signature} transfers the signature.
\end{proof}

The important logical point is that the metric is \emph{derived after}
regular soldering.  It is not a field of the base-gluing datum and is not
inserted into the primitive local data.

\subsection{Orientation reduction}

Let
\[
   F_{ij}(y)
   :=
   A_j(\phi_{ij}(y))^{-1}
   \circ T_{ij}(y)\circ A_i(y).
\tag{8.3}
\]
By the solder law this is the appropriate projected internal Lorentz
transition (with the inverse-index convention induced by
\eqref{eq:regular-solder-law}).  Therefore
\[
   \det F_{ij}(y)=1.
\tag{8.4}
\]

\begin{proposition}[Project-native orientation reduction]
\label{prop:orientation-reduction}
The regular solder frames form a coherently oriented family: their frame
changes have determinant \(1\).  Equivalently, the determinant cocycle of
the genuine tangent transitions is a continuous coboundary determined by
the solder-frame determinants.
\end{proposition}

\subsection{Future-cone reduction}

Define the intrinsic interior future cone exactly by
\[
   \mathcal C^{+}_{\mathrm{int}}
   :=
   \{x\in\Smodel : x\in C_{\Smodel}\ \text{and}\ N(x)>0\}.
\]
This uses only the already constructed square cone and intrinsic quadratic
form.  In chart \(i\), define
\[
   \mathcal C^{+}_{i}(y)
   :=
   A_i(y)\mathcal C^{+}_{\mathrm{int}}.
\tag{8.5}
\]

\begin{proposition}[Project-native future-cone reduction]
\label{prop:future-cone}
For \(y\in W_{ij}\),
\[
   u\in\mathcal C^{+}_{i}(y)
   \quad\Longleftrightarrow\quad
   T_{ij}(y)u
   \in
   \mathcal C^{+}_{j}(\phi_{ij}(y)).
\tag{8.6}\label{eq:future-cone-transform}
\]
Thus the regular solder determines a coherent chart-independent field of
future cones on \(TM\).
\end{proposition}

The local vector field
\[
   \tau_i(y):=A_i(y)\mathbf1_{\Smodel}
\tag{8.7}
\]
is smooth, future-directed, and unit timelike.  It need not glue to a
single global field.  More precisely, the distinguished fields
\(\tau_i\) glue across an overlap exactly when the corresponding projected
Lorentz transition fixes \(\mathbf1_{\Smodel}\).

\begin{boundary}[Time-orientability terminology]
\label{bdry:timeorientation}
The conclusion formally established here is the coherent future-cone
reduction \eqref{eq:future-cone-transform}.  No comparison theorem in the present development
identifies this project-native structure with every standard equivalent
definition of time orientability.  Likewise, existence of an arbitrary
global smooth future-directed timelike section would require an additional
global section theorem in the present bundle packaging.
\end{boundary}

% ======================================================================
\clearpage
\section{Native \Cech\ lifting and global reconvergence}
\label{sec:reconvergence}
% ======================================================================

We now compare the bottom-up construction with the obstruction-theoretic
language that ordinarily appears in Spin geometry.

\subsection{Fixed-cover lifting obstruction}

Let
\[
   g_{ij}:U_i\cap U_j\to\LorG
\]
be a visible Lorentz cocycle.  A family of local Spin lifts consists of
continuous maps
\[
   h_{ij}:U_i\cap U_j\to\SpinG,
   \qquad
   \rho(h_{ij})=g_{ij}.
\tag{9.1}
\]
Local sections of \(\rho\) make such lifts available locally.  On triple
overlaps, define the defect
\[
   \epsilon_{ijk}
   :=
   h_{ij}h_{jk}h_{ik}^{-1}.
\tag{9.2}
\]
Because \(\rho(\epsilon_{ijk})=1\),
\[
   \epsilon_{ijk}\in\ker\rho=\{\pm1\}.
\tag{9.3}
\]
This defines a fixed-cover kernel-valued obstruction class.
The independence from the chosen lift family is elementary at this level.
If \(h'_{ij}\) is another lift of the same visible Lorentz cocycle, then
there are kernel-valued functions
\(\alpha_{ij}:U_i\cap U_j\to\{\pm1\}\) with
\[
    h'_{ij}=\alpha_{ij}h_{ij}.
\]
Because the kernel is central,
\[
\begin{aligned}
 \epsilon'_{ijk}
 &=h'_{ij}h'_{jk}(h'_{ik})^{-1}\\
 &=\alpha_{ij}\alpha_{jk}\alpha_{ik}^{-1}\,
   \epsilon_{ijk}.
\end{aligned}
\tag{9.3a}\label{eq:obstruction-coboundary}
\]
Thus two representatives differ by the fixed-cover kernel-valued
coboundary \(\delta\alpha\), and determine the same project-native class.
A coherent native Spin cocycle gives \(\epsilon_{ijk}=1\) identically, so
its class is trivial.

This is the familiar mechanism underlying the usual obstruction theory for
Spin structures \cite{Steenrod,MilnorStasheff,LawsonMichelsohn}; however,
the present class is kept at its exact project-native fixed-cover level.

\begin{boundary}[No Stiefel--Whitney identification]
\label{bdry:no-sw}
No theorem here identifies the determinant obstruction with \(w_1(TM)\),
or the fixed-cover kernel-valued lifting obstruction with \(w_2(TM)\).
No refinement limit or comparison with global
\(H^{1}(M;\Ztwo)\) or \(H^{2}(M;\Ztwo)\) is used.  Consequently the
statement \(w_1(TM)=w_2(TM)=0\) is not a theorem of the present paper.
\end{boundary}

\subsection{Reconvergence under the solder gate}

The regular solder turns the projected internal Lorentz cocycle into the
actual tangent Lorentz cocycle read in the solder frames.

\begin{theorem}[Solder-induced cocycle identification]
\label{thm:cocycle-identification}
For every regular solder and every overlap,
the Lorentz transformation obtained by comparing tangent components in the
two solder frames is exactly the projected native transition \(g_{ij}\).
\end{theorem}

This is a rearrangement of the intertwining equation \eqref{eq:regular-solder-law}, but it is
the central bridge: it allows the internal double-cover information to be
read as information about the tangent bundle.

\begin{theorem}[Global reconvergence]
\label{thm:reconvergence}
Assume \(B\) is smoothly glued and a regular solder exists for a native Spin
cocycle \(\widetilde g\).  Then:
\begin{enumerate}
\item there are continuous nowhere-zero functions
      \(a_i:D_i\to\R\) such that
\[
  \det T_{ij}(y)\,a_i(y)
  =
  a_j(\phi_{ij}(y));
\tag{9.4}\label{eq:orientation-coboundary}
\]
\item the tangent Lorentz cocycle in solder frames equals
      \(\rho(\widetilde g_{ij})\);
\item every local lift family of that projected cocycle has trivial
      project-native fixed-cover Spin-lifting obstruction;
\item the project-native future-cone reduction of
      Proposition~\ref{prop:future-cone} exists.
\end{enumerate}
\end{theorem}

\begin{proof}
For (1), take the determinant of the regular-solder intertwining law.
The internal Lorentz factor has determinant one, so the determinant of the
tangent transition is the ratio of the solder-frame determinants, giving
\eqref{eq:orientation-coboundary}.  Statement (2) is
\cref{thm:cocycle-identification}.  For (3), the native Spin cocycle itself
is a coherent lift of its projection; the triple-overlap defect is
therefore trivial.  Independence of the obstruction from the chosen local
lift family transfers triviality to every lift family.  Statement (4) is
Proposition~\ref{prop:future-cone}.
\end{proof}

Two useful contrapositives are immediate.

\begin{corollary}[Orientation gate]
\label{cor:orientation-gate}
If no continuous nowhere-zero \(0\)-cochain \(a_i\) can satisfy
\eqref{eq:orientation-coboundary}, then no regular solder exists for any native Spin datum over
that base gluing.
\end{corollary}

\begin{corollary}[Spin-lifting gate]
\label{cor:spin-gate}
If a visible Lorentz cocycle has nontrivial project-native lifting
obstruction, then it admits no coherent native Spin transition datum
projecting to it.  In particular it cannot arise as the tangent Lorentz
cocycle of a regularly soldered native Spin datum.
\end{corollary}

The converse of \cref{thm:reconvergence} is not proved.  Vanishing of these
two necessary obstructions is therefore not asserted to be sufficient for
a regular solder.

% ======================================================================
\clearpage
\section{Controls and residual global Spin freedom}
\label{sec:controls}
% ======================================================================

The previous theorems are accompanied by controls designed to distinguish
logical necessity from convenient examples.

\subsection{One-chart positive control}

On a one-chart base there are no nontrivial overlap constraints.  The base
is smooth, a regular solder exists, and there is no artificial sign family
generated by the cover.  This verifies that the closure machinery does not
manufacture an obstruction where none is present.

\subsection{Orientation-reversing negative control}

There is a smooth M\"obius-type base gluing whose cross-overlap has two
relevant components.  On the wrap-around component the genuine tangent
transition has determinant \(-1\), whereas on the ordinary overlap component
it has determinant \(+1\).

\begin{control}[Orientation reversal excludes regular solder]
\label{control:moebius}
For the orientation-reversing control, no regular solder exists for
\emph{any} native Spin transition datum.
\end{control}

\begin{proof}[Reason]
Every projected native Spin transition lies in \(\LorG\) and hence has
determinant \(+1\).  If a regular solder existed, the determinant of each
local solder frame would be a continuous nowhere-zero function on a
connected chart domain and therefore would have constant sign there.  The
solder determinant law on the wrap-around component, where the tangent
determinant is \(-1\), forces opposite signs in the two charts; the same law
on the ordinary component, where the tangent determinant is \(+1\), forces
the same signs.  These requirements are incompatible.  This is the concrete
mechanism behind the orientation gate Corollary~\ref{cor:orientation-gate}; the
presence of both overlap components is essential.
\end{proof}

Native Spin transition data nevertheless exist over this same gluing.
Thus existence of Spin-valued fibre data is strictly weaker than existence
of a Lorentz solder coupling those data to \(TM\).

\subsection{Combinatorial-label negative control}

Consider also a deliberately abstract Spin-foam-style label structure.  For arbitrary such combinatorial labels,
the same label input is compatible with the one-chart positive control,
where a regular solder exists, and with the M\"obius-type control, where no
regular solder exists for any native Spin datum.  Therefore these
combinatorial labels do not determine a regular solder.  This is a statement
about this abstract control structure, not about any published spin-foam
model or about spin foams in general.

\subsection{Fixed-cover kernel freedom}

Because \(\ker\rho=\{\pm1\}\), a native lift can retain discrete global
freedom even when the projected Lorentz transition system is fixed.

\begin{proposition}[One-loop sign freedom]
\label{prop:one-loop-freedom}
On the fixed one-loop control cover there are two native Spin transition
systems that:
\begin{enumerate}
\item have the same projected Lorentz transition system;
\item both admit regular solders;
\item are not equivalent under the project's kernel-valued fixed-cover
      gauge relation.
\end{enumerate}
\end{proposition}

\begin{proposition}[Finite fixed-cover multiplicity]
\label{prop:multi-loop}
For \(N\) independent loop labels, the construction gives an injection of
\(\{\pm1\}^{N}\) into the fixed-cover gauge classes.  In particular, for
\(N=2\) there are four pairwise distinguished classes, all compatible with
regular soldering in the control.
\end{proposition}

\begin{proof}[Proof of Propositions~\ref{prop:one-loop-freedom} and~\ref{prop:multi-loop}]
For each loop \(n\), the cross-overlap of its two patches has two connected
components: the wrap-around component and the ordinary component.  Given a
sign assignment \(\sigma\in\{\pm1\}^{N}\), define a kernel-valued
\Cech~1-cocycle \(z^\sigma\) which equals \(\sigma_n\) on the
wrap-around component of loop \(n\) and equals \(1\) on the ordinary
component and on all unaffected overlaps.  Since every value of
\(z^\sigma\) lies in \(\ker\rho=\{\pm1\}\), twisting the trivial Spin
cocycle by \(z^\sigma\) leaves the projected Lorentz cocycle unchanged.
Consequently the same regular solder for the periodic control works for every
sign assignment.

It remains to prove that distinct sign assignments are not related by a
kernel-valued fixed-cover gauge.  Suppose \(\sigma_n\ne\tau_n\) for some
loop \(n\), and assume that patchwise kernel-valued gauge functions
\(\lambda_-\) and \(\lambda_+\) relate the corresponding twisted
cocycles.  The kernel is discrete and each patch is connected, so continuity
forces \(\lambda_-\) and \(\lambda_+\) to be constant on their patches.
Evaluating the gauge law at a point of the ordinary overlap, where both sign
cocycles equal \(1\), gives
\[
    \lambda_-\lambda_+^{-1}=1.
\]
Evaluating the same law at a point of the wrap-around component gives, using
centrality of the kernel,
\[
    \tau_n
      =\lambda_-\,\sigma_n\,\lambda_+^{-1}
      =\sigma_n,
\]
contradicting \(\sigma_n\ne\tau_n\).  Hence different sign assignments
give different fixed-cover gauge classes.  For one loop this yields exactly
the two classes of Proposition~\ref{prop:one-loop-freedom}; for \(N\)
independent loops it gives the asserted injection, and for \(N=2\) the four
sign assignments give four pairwise inequivalent classes.
\end{proof}

\begin{boundary}
The preceding statements do not compute \(H^{1}(M;\Ztwo)\).  They are
fixed-cover statements; no refinement argument identifying them with
global cohomology is part of the present construction.
\end{boundary}

This control is important conceptually: bottom-up Spin closure does not
canonize the Spin lift.

% ======================================================================
\clearpage
\section{A selected one-parameter transition deformation}
\label{sec:deformation}
% ======================================================================

The static closure structure is now complete enough to perform a first
deformation smoke test.  The deformation is placed at the strongest
transition-level object that presently exists: the native Spin cocycle.
There is no connection one-form, path-ordered transport, holonomy
functional, curvature two-form, or field equation in the construction.

\subsection{Selected intrinsic one-parameter subgroup}

Choose the first Euclidean basis direction \(\varepsilon_{0}\in V\)
and its Clifford image
\[
    \gamma_{0}:=\iota(\varepsilon_{0}),
    \qquad \gamma_{0}^{2}=1.
\]
The generator \(\gamma_{0}\) is selected for the control; no theorem makes
this direction canonical.

\begin{definition}[Selected Spin path]
\label{def:G-lambda}
For \(\lambda\in\R\), define
\[
   G(\lambda)
   :=
   \cosh\!\left(\frac{\lambda}{2}\right)
   +
   \sinh\!\left(\frac{\lambda}{2}\right)\gamma_{0}
   \in\Cl_{3}.
\tag{11.1}\label{eq:selected-spin-path}
\]
\end{definition}

The hyperbolic identities and \(\gamma_{0}^{2}=1\) give
\[
   G(\lambda)c(G(\lambda))=1,
\]
hence \(G(\lambda)\in\SpinG\).  Moreover,
\[
   G(0)=1,\qquad
   G(\lambda+\mu)=G(\lambda)G(\mu),\qquad
   G(-\lambda)=G(\lambda)^{-1}.
\tag{11.2}
\]
The parameterization is injective.

The projection \(\rho(G(\lambda))\) is nontrivial for
\(\lambda\neq0\).  On the plane spanned by the intrinsic scalar direction
and \((0,\varepsilon_{0})\), its action is the standard boost block
\[
   \begin{pmatrix}
      \cosh\lambda & \sinh\lambda\\
      \sinh\lambda & \cosh\lambda
   \end{pmatrix},
\tag{11.3}\label{eq:boost-block}
\]
and it acts trivially on the two orthogonal Euclidean directions.
In particular,
\[
   \rho(G(\lambda))\mathbf1_{\Smodel}
   =
   \bigl(\cosh\lambda,\sinh\lambda\,\varepsilon_{0}\bigr).
\tag{11.4}
\]

\begin{boundary}[Selected control slice]
\label{bdry:selected}
The family \eqref{eq:selected-spin-path} is a deliberately selected one-parameter subgroup
generated by one chosen Clifford direction.  It is neither proved unique
nor canonical, and it does not exhaust the possible intrinsic Spin
deformation directions.
\end{boundary}

\subsection{The fixed periodic control}

Fix the periodic base gluing used by the control, together with its fixed
cover and two-component periodic overlap.  We make no identification of
this quotient with \(S^{1}\times\R^{3}\), and no fundamental-group theorem
is required.

The native Spin transition cocycle is deformed using \(G(\lambda)\).
The base gluing and hence the tangent transitions are held fixed; the
projected Lorentz transition is always obtained by applying \(\rho\) to
the Spin transition.  The regular solder is then the variable to be solved
for.

Let
\[
  \cR_A
  :=
  \left\{
      \lambda\in\R:
      \text{a regular solder exists for the deformed cocycle at }\lambda
  \right\}.
\tag{11.5}
\]

\begin{theorem}[Admissibility locus]
\label{thm:admissibility}
For the fixed periodic control,
\[
   \cR_A=\R.
\tag{11.6}\label{eq:admissibility-region}
\]
Thus every finite parameter value admits an explicitly constructed regular
solder.  In particular admissibility is symmetric under
\(\lambda\mapsto-\lambda\).
\end{theorem}

\begin{proof}
Let \(\chi:\R\to[0,1]\) be a fixed \(C^\infty\) transition profile with
\(\chi(t)=0\) for \(t\le3\) and \(\chi(t)=1\) for \(t\ge4\).  On the two
patches of each periodic control loop define
\[
  \alpha_\lambda(p,y)=
  \begin{cases}
     0, & p\text{ is the first patch},\\
     \lambda\,\chi(y_0), & p\text{ is the second patch}.
  \end{cases}
\tag{11.6a}\label{eq:solder-angle}
\]
Let \(P(t)=\rho(G(t))\in\LorG\).  By the one-parameter subgroup law,
\(P(s+t)=P(s)P(t)\) and \(P(-t)=P(t)^{-1}\), and the closed form
\eqref{eq:boost-block} makes \(t\mapsto P(t)\) smooth.  Define
\[
   A_p^{(\lambda)}(y):=P(\alpha_\lambda(p,y)).
\tag{11.6b}\label{eq:explicit-solder}
\]
This is the identity on the first patch and smoothly interpolates on the
second patch between the identity and \(P(\lambda)\).

All tangent transitions of the frozen periodic control are the identity.
The projected deformed Spin transition on an overlap is
\(P(\beta_{pq}^{(\lambda)}(x))\), where
\(\beta_{pq}^{(\lambda)}\) is \(\pm\lambda\) on the wrap-around component
and \(0\) on the ordinary component.  By the defining values of \(\chi\),
one verifies on each overlap component that
\[
   \alpha_\lambda(q,\phi_{pq}(y))
   =\alpha_\lambda(p,y)+\beta_{pq}^{(\lambda)}(x_y).
\tag{11.6c}\label{eq:solder-angle-step}
\]
Therefore
\[
\begin{aligned}
  A_q^{(\lambda)}(\phi_{pq}(y))
   &=P(\alpha_\lambda(q,\phi_{pq}(y)))\\
   &=P(\alpha_\lambda(p,y))
     P(\beta_{pq}^{(\lambda)}(x_y))\\
   &=A_p^{(\lambda)}(y)\,g_{pq}^{(\lambda)}(x_y),
\end{aligned}
\]
which is precisely the regular-solder law because the tangent transition is
\(\id\).  Smoothness and pointwise invertibility follow from those of
\(P\) and \(\chi\).  Hence a regular solder exists for every
\(\lambda\in\R\), proving \(\cR_A=\R\).
\end{proof}

The result is not vacuous: the project also contains closure families for
which not every parameter is admissible.  Equation \eqref{eq:admissibility-region} is
therefore a property of this particular control family, not a consequence
of the definition of admissibility.

\subsection{Full Spin \Cech\ gauge orbit}

Let \(\widetilde g^{(\lambda)}\) denote the native transition cocycle at
parameter \(\lambda\).

\begin{theorem}[Gauge classification of the selected family]
\label{thm:gauge-orbit}
For all \(\lambda_{1},\lambda_{2}\in\R\),
\[
   \widetilde g^{(\lambda_{1})}
   \sim_{\mathrm{full\ Spin\ \Cech}}
   \widetilde g^{(\lambda_{2})}.
\tag{11.7}
\]
The Spin-side gauge is continuous.  After projection by \(\rho\), its
Lorentz representative is \(C^\infty\) on every chart domain, and the
projected Lorentz cocycles also lie in one gauge orbit.
\end{theorem}

\begin{proof}
First compare an arbitrary \(\lambda\) with the neutral value.  Using the
same interpolation function \(\alpha_\lambda\) as in
\eqref{eq:solder-angle}, define on patch \(p\)
\[
   u_p^{(\lambda)}(x)
   :=G\!\left(-\alpha_\lambda
       \bigl(p,\operatorname{coord}_p(x)\bigr)\right).
\tag{11.7a}\label{eq:spin-gauge}
\]
The map is continuous because \(G(t)\) is continuous in the intrinsic Spin
topology and the chart-coordinate profile is continuous.  On an overlap,
\eqref{eq:solder-angle-step} and the group law of \(G\) give
\[
   \widetilde g_{pq}^{(\lambda)}(x)
   =u_p^{(\lambda)}(x)\,
      \widetilde g_{pq}^{(0)}(x)\,
      \bigl(u_q^{(\lambda)}(x)\bigr)^{-1}.
\tag{11.7b}\label{eq:spin-coboundary}
\]
Thus every parameter value is full-Spin-\Cech\ gauge equivalent to zero.
For two arbitrary values \(\lambda_1,\lambda_2\), compose the inverse gauge
from \(\lambda_1\) to zero with the gauge from zero to \(\lambda_2\).

Applying \(\rho\) to \eqref{eq:spin-coboundary} gives the corresponding
Lorentz gauge relation.  Its chartwise representative is
\(P(-\alpha_\lambda)\); by the explicit hyperbolic formula
\eqref{eq:boost-block} and smoothness of \(\alpha_\lambda\), this projected
gauge is \(C^\infty\) on every chart domain.  No smooth structure on the
intrinsic Spin group is used or claimed.
\end{proof}

The narrower kernel-valued fixed-cover equivalence behaves differently.
For \(\lambda\neq0\), the deformed datum is not merely a \(\{\pm1\}\)-lift
relabeling of the neutral datum, because its Lorentz projection itself
moves.  There is no contradiction: the two equivalence relations answer
different questions.

\subsection{All regular solder solutions, not only one branch}

For a parameter \(\lambda\), write
\[
   \cE_{\lambda}
   :=
   \left\{
      \text{regular solders for }
      (B_{\mathrm{per}},\widetilde g^{(\lambda)})
   \right\}.
\tag{11.8}
\]

\begin{theorem}[Solution-space correspondence]
\label{thm:solution-correspondence}
For every \(\lambda_{1},\lambda_{2}\in\R\), there is an explicit bijection
\[
   \mathcal T_{\lambda_{1},\lambda_{2}}:
   \cE_{\lambda_{1}}
   \xrightarrow{\;\sim\;}
   \cE_{\lambda_{2}},
\tag{11.9}\label{eq:solution-bijection}
\]
obtained by precomposing the local solder comparison with the projected
gauge relating the two native Spin cocycles.  The inverse is
\(\mathcal T_{\lambda_{2},\lambda_{1}}\).
\end{theorem}

\begin{proof}
Set
\[
   \delta_{12}(p,y)
   :=\alpha_{\lambda_2}(p,y)-\alpha_{\lambda_1}(p,y),
   \qquad
   Q_{12,p}(y):=P(\delta_{12}(p,y)).
\tag{11.9a}\label{eq:gauge-angle}
\]
For a regular solder \(A\in\cE_{\lambda_1}\), define
\[
   \bigl(\mathcal T_{\lambda_1,\lambda_2}A\bigr)_p(y)
   :=A_p(y)\circ Q_{12,p}(y).
\tag{11.9b}\label{eq:solder-transport}
\]
The fields \(Q_{12,p}\) are smooth and invertible.  Subtracting the two
identities \eqref{eq:solder-angle-step} for \(\lambda_1\) and
\(\lambda_2\) gives exactly the overlap equation needed to convert the
\(\lambda_1\)-solder law into the \(\lambda_2\)-solder law.  Thus
\eqref{eq:solder-transport} is a regular solder at \(\lambda_2\).

Since
\[
   \delta_{21}=-\delta_{12}
   \quad\text{and}\quad
   P(-t)=P(t)^{-1},
\]
transporting from \(\lambda_1\) to \(\lambda_2\) and then back composes
with \(Q_{12}Q_{21}=\id\) pointwise.  Hence
\(\mathcal T_{\lambda_2,\lambda_1}\) is the inverse of
\(\mathcal T_{\lambda_1,\lambda_2}\), proving the stated bijection.
\end{proof}

\begin{theorem}[Metric preservation]
\label{thm:metric-preservation}
Let \(A\in\cE_{\lambda_{1}}\), and let
\(A'=\mathcal T_{\lambda_{1},\lambda_{2}}(A)\).
Then for every \(x\in M\) and \(X,Y\in T_xM\),
\[
   g_{A'}(X,Y)=g_A(X,Y).
\tag{11.10}
\]
\end{theorem}

\begin{proof}
In a chart, \eqref{eq:solder-transport} gives
\[
    A'_p=A_p\circ Q_{12,p},
    \qquad
    (A'_p)^{-1}=Q_{12,p}^{-1}\circ A_p^{-1}.
\]
Therefore
\[
\begin{aligned}
 g_{A',p}(u,v)
 &=B\!\left(Q_{12,p}^{-1}A_p^{-1}u,
            Q_{12,p}^{-1}A_p^{-1}v\right)\\
 &=B\!\left(A_p^{-1}u,A_p^{-1}v\right)\\
 &=g_{A,p}(u,v),
\end{aligned}
\]
because \(Q_{12,p}(y)=P(\delta_{12}(p,y))\in\LorG\) preserves \(B\).
The equality holds in every chart and is compatible with the tangent
transition law, hence it is the equality of the induced tangent metrics at
every base point and on every pair of tangent vectors.
\end{proof}

\begin{corollary}[Scoped negative smoke-test conclusion]
\label{cor:smoketest}
On the fixed periodic control, the selected transition-cocycle parameter
has no closure-selection power and does not change the induced tangent
Lorentz metric under the certified correspondence of regular-solder
solutions.
\end{corollary}

\begin{boundary}[What Corollary~\ref{cor:smoketest} does not say]
The result does not classify arbitrary native Spin deformations, arbitrary
covers, co-varying bases, or connection-level deformations.  It does not
show that ``all transition deformations are pure gauge'', and it carries no
curvature interpretation.  The parameter \(\lambda\) is the coordinate of
the selected Spin subgroup \eqref{eq:selected-spin-path}, not a curvature scalar.
\end{boundary}

% ======================================================================
\clearpage
\section{Mathematical boundary and next structural edge}
\label{sec:boundary}
% ======================================================================

The construction reaches a sharp stopping point.  It is useful to state it
positively before listing what remains open.

From the sole local primitive \(V=\R^{3}\) with its Euclidean form, one has
constructed the Lorentz spin factor, the Euclidean Clifford envelope, the
intrinsic Spin double cover, and the six-dimensional Lorentz-skew algebra.
After adding explicit base-gluing data and sufficient topological and smooth
conditions, one obtains a smooth four-manifold.  After adding a regular
solder, one obtains a smooth Lorentzian tangent metric, coherent
orientation and future-cone reductions, and the exact identification of
the tangent Lorentz cocycle with the projection of native Spin transition
data.

What is \emph{not} yet part of the mathematical object is equally precise.

\begin{enumerate}
\item The intrinsic topological group \(\SpinG\) has not been packaged and
      related to \(\gB\) as a Lie group/Lie algebra pair:
\[
   \operatorname{Lie}(\SpinG)\simeq\gB
   \qquad\text{is open.}
\tag{12.1}
\]
\item The differential \(d\rho_e\) of the native Spin cover has not been
      constructed.
\item No Spin or Lorentz connection one-form has been defined.
\item Consequently there is no covariant derivative, parallel transport,
      holonomy, curvature, torsion, action functional, or field equation in
      the present theory.
\item The fixed-cover obstruction classes have not been compared
      formally with \(w_1(TM)\) and \(w_2(TM)\).
\item The project-native future-cone reduction has not been identified by
      theorem with all standard formulations of time orientability.
\end{enumerate}

The selected deformation result of \cref{sec:deformation} gives a concrete
reason not to skip these missing edges.  A nontrivial curve of
transition-level Spin data can move the projected Lorentz cocycle while
remaining fully gauge absorbable on the tested control and leaving the
tangent metric invariant.  Thus, on this control, the first tested transition-level knob has no
closure-selection power and does not change the induced tangent Lorentz
metric under the certified solder correspondence.  This does not prove that a
connection-level theory will be nontrivial; it identifies the next
mathematically meaningful place at which such a question can be posed.

% ======================================================================
\clearpage
\section{Conclusion}
% ======================================================================

The main structural lesson is a separation theorem for levels of geometric
information.

At the local level, positive Euclidean three-space already contains enough
algebraic structure to produce a four-dimensional Lorentz spin factor, its
Clifford envelope, an intrinsic Spin double cover of an intrinsic Lorentz
group, and a six-dimensional Lorentz-skew Lie algebra.  This part is
intrinsic to the local input.

At the global level, incidence is independent information.  Local pieces,
even when equipped with fibre transitions, do not determine which points
of different pieces form a common base.  Explicit base gluing is therefore
a genuine additional primitive.  Under a closed-graph separation condition,
countability, and smooth compatibility, it reconstructs a smooth
four-manifold.

At the tangent level, a second independent gap appears.  The derivative
cocycle \(D\phi_{ij}\) of that manifold is not forced to equal the internal
Lorentz cocycle.  A regular solder is the additional gate that identifies
them in the correct bundle-theoretic sense.  Once supplied, it yields a
smooth Lorentzian tangent metric and forces the project-native orientation,
future-cone, and Spin-lifting compatibility conditions.

The resulting implication is therefore conditional and exact:
\[
\begin{array}{c}
(\R^{3},\langle\cdot,\cdot\rangle)
\Longrightarrow
(\Smodel,N,\Cl_{3},\SpinG\!\to\!\LorG,\gB)
\\[0.5em]
\begin{gathered}
\text{local core}+\text{base gluing}\\
+\text{closed gluing graph}+\text{countable index}+\text{smoothness}
\end{gathered}
\Longrightarrow M^{4}
\\[0.5em]
M^{4}+\text{native Spin cocycle}+\text{regular solder}
\\[-0.1em]
\Downarrow
\\[-0.1em]
(TM,g),
\\
\text{project-native orientation compatibility},
\\
\text{project-native future-cone reduction},
\\
\text{tangent Lorentz cocycle = projected native Spin cocycle},
\\
\text{trivial project-native fixed-cover Spin-lifting obstruction}.
\end{array}
\tag{13.1}\label{eq:final-chain}
\]
No arrow in \eqref{eq:final-chain} hides a missing global existence theorem.

The first deformation smoke test is correspondingly modest but sharp.  On
the selected fixed periodic control, the chosen one-parameter Spin
transition family is globally admissible and gauge absorbable for all
finite parameters, and its entire regular-solder solution spaces are
identified without changing the induced tangent Lorentz metric.  This
closes that specific control question and leaves the infinitesimal
Spin-to-Lorentz bridge, and only thereafter connection-level geometry, as
the next structural problem.

% ======================================================================
\appendix
\clearpage
\section{Dependency and status table}
\label{app:status}
% ======================================================================

\begin{center}
\small
\begin{tabular}{>{\raggedright\arraybackslash}p{0.30\textwidth}
                >{\raggedright\arraybackslash}p{0.23\textwidth}
                >{\raggedright\arraybackslash}p{0.37\textwidth}}
\toprule
Object or statement & Status & Dependency \\
\midrule
\(V=\R^{3}\), Euclidean form & primitive input & none \\
\(\Smodel,N,B\), signature \((1,3)\) & \statusderived & Euclidean input \\
\(\Cl_{3}\), paravectors & \statusderived & \(q_{3}\) \\
\(\SpinG\to\LorG\), kernel \(\{\pm1\}\) & \statusderived & Clifford core \\
\(\gB\simeq\Lambda^{2}\Smodel\) & \statusderived & \(B\) \\
Base incidence/gluing & \statusdatum & not determined by local pieces \\
Closed gluing graph & additional hypothesis & needed for certified Hausdorffness \\
Countable index set & additional hypothesis & needed for certified second countability \\
Smooth gluing & sufficient additional hypothesis & base transitions \(\phi_{ij}\) \\
Smooth \(M^{4}\) & \statusderived & preceding global hypotheses \\
Tangent transitions \(D\phi_{ij}\) & \statusderived & smooth atlas \\
Native Spin cocycle & global datum & fixed cover of \(M\) \\
Regular solder & \statusdatum & not generally inhabited \\
Smooth Lorentz tangent metric & \statusderived & regular solder \\
Orientation reduction & \statusderived & regular solder \\
Future-cone reduction & \statusderived & regular solder \\
Project-native Spin obstruction trivial & \statusderived & regular solder + native cocycle \\
\(w_1,w_2\) comparison & \statusopen & characteristic-class comparison theorem \\
\(\operatorname{Lie}(\SpinG)\simeq\gB\) & \statusopen & smooth intrinsic Spin theory \\
\(d\rho_e\) & \statusopen & preceding infinitesimal bridge \\
Connection, holonomy, curvature & \statusopen & connection-level extension \\
\bottomrule
\end{tabular}
\end{center}

% ======================================================================
\clearpage
\section{Formal source map}
\label{app:source-map}
% ======================================================================

This appendix is a reproducibility map, not part of the mathematical
dependency argument.  The human-readable theorem statements above are tied
to the current formal source by the following principal declarations.

\begin{center}
\footnotesize
\begin{tabular}{>{\raggedright\arraybackslash}p{0.28\textwidth}
                >{\raggedright\arraybackslash}p{0.64\textwidth}}
\toprule
Paper statement & Principal formal endpoint(s) \\
\midrule
Spin factor and Lorentz form &
\path{SpinCore.sJ_jordan},
\path{SpinCore.lorentz_quadratic_form} \\
Jordan--Clifford bridge &
\path{SpinCore.clifford_jordan},
\path{SpinCore.adjoin_spinToCl_eq_top} \\
Intrinsic double cover &
\path{SpinCore.spin_double_cover_bundled},
\path{SpinCore.spin_double_cover_topological_endpoint},
\path{SpinCore.spin_double_cover_local_section_endpoint} \\
Bivector/Lorentz-skew algebra &
\path{SpinCore.bivectorEquivSkew},
\path{SpinCore.finrank_gB},
\path{SpinCore.gBLie_bracket} \\
Base non-derivability &
\path{EmergentBase.base_not_determined_by_local_pieces} \\
Topological reconstruction &
\path{EmergentBase.BaseGluingData.emergentManifoldCertificate} \\
Smooth manifold &
\path{EmergentBase.BaseGluingData.isManifold_of_smoothGluing},
\path{EmergentBase.BaseGluingData.smoothEmergentManifoldCertificate} \\
Genuine tangent transition &
\path{EmergentBase.BaseGluingData.tangentTransition_eq_derivative_baseTransition} \\
Tangent/internal independence &
\path{SpinNative.tangent_spin_base_data_do_not_force_transition_identification} \\
Regular solder &
\path{SpinNative.SmoothTangentSolderData};
\path{SpinNative.smooth_solder_iff_regular_bundle_equivalence} \\
Smooth tangent metric &
\path{SpinNative.SmoothTangentSolderData.smooth_tangent_lorentz_metric} \\
Orientation/future cone &
\path{SpinNative.SmoothTangentSolderData.orientation_time_orientation_status} \\
Global reconvergence &
\path{Task36.classical_reconvergence},
\path{Task36.spin_obstruction_nonzero_implies_no_solder} \\
Adversarial controls &
\path{Task36.adversarial_global_topology_certificate},
\path{Task36.spin_foam_labels_do_not_determine_solder} \\
Selected deformation path &
\path{Task37.Deformation.transportState_add},
\path{Task37.Deformation.projectedTransportState_apply} \\
Fixed-control admissibility &
\path{Task37.Deformation.loopParaSolder} (construction),
\path{Task37.Deformation.controlAAdmissible_all},
\path{Task37.Deformation.regularRegion_loopTransportFamily} \\
Full Spin \Cech\ gauge orbit &
\path{Task37.Deformation.loopGaugeFun} (construction),
\path{Task37.Deformation.spinCocycle_gaugeEquiv},
\path{Task37.Deformation.projectedCocycle_gaugeEquiv} \\
Selected deformation smoke test &
\path{Task37.Deformation.native_transport_knob_smoketest_certificate} \\
Whole solution-space transport &
\path{Task37.Deformation.loopSolderEquiv} (construction),
\path{Task37.Deformation.loopSolderTransport_involutive},
\path{Task37.Deformation.loopSolderSolutionSpace_correspondence},
\path{Task37.Deformation.tangentMetric_loopSolderTransport} \\
\bottomrule
\end{tabular}
\end{center}

The formal development is available at
\url{https://github.com/antaris82/The_Relational_Spin-Closure_Hypothesis}
\cite{RSCHFormal}.
The main development is written in Lean~4 and uses Mathlib
\cite{LeanSystem,Mathlib2020}.  The mathematical claims of this paper should
be read at the strength of the displayed hypotheses above; the formal
artifact supplies machine-checked proofs of the corresponding registered
source statements.

% ======================================================================
\clearpage
\section*{Acknowledgments}
\addcontentsline{toc}{section}{Acknowledgments}
% ======================================================================

The editor thanks Thomas Stoer and Thomas Klimpel for sustained, critical,
and constructive discussions.  These exchanges repeatedly helped clarify
assumptions, expose weak formulations, and sharpen the mathematical
questions addressed in this work.

% ======================================================================
\section*{AI and formalization disclosure}
\addcontentsline{toc}{section}{AI and formalization disclosure}
% ======================================================================

ChatGPT (OpenAI), acting as AI Assistant, was used for mathematical
exploration, source-level analysis, organization, consistency checking, and
editorial drafting.  Aristotle was used as Lean Agent for prover-assisted
formalization and verification.  The editor curated the scope and the
presentation.  Mathematical claims are asserted only at the strength stated
in the paper and are tied, where applicable, to the referenced formal Lean
endpoints.  The formal artifact, rather than the language-model prose, is
the machine-checked source for the certified statements.

% ======================================================================
\small
\setlength{\emergencystretch}{2em}
\begin{thebibliography}{99}

\bibitem{Chevalley}
C.~Chevalley,
\emph{The Algebraic Theory of Spinors},
Columbia University Press, New York, 1954.

\bibitem{FarautKoranyi}
J.~Faraut and A.~Kor\'anyi,
\emph{Analysis on Symmetric Cones},
Oxford Mathematical Monographs,
Oxford University Press, Oxford, 1994.
doi:10.1093/oso/9780198534778.001.0001.

\bibitem{Husemoller}
D.~Husemoller,
\emph{Fibre Bundles},
3rd ed., Graduate Texts in Mathematics 20,
Springer, New York, 1994.

\bibitem{JordanNeumannWigner}
P.~Jordan, J.~von Neumann, and E.~Wigner,
On an algebraic generalization of the quantum mechanical formalism,
\emph{Annals of Mathematics} \textbf{35} (1934), 29--64.
doi:10.2307/1968117.

\bibitem{KobayashiNomizu}
S.~Kobayashi and K.~Nomizu,
\emph{Foundations of Differential Geometry, Vol.~I},
Interscience Publishers, New York, 1963.

\bibitem{LawsonMichelsohn}
H.~B.~Lawson, Jr. and M.-L.~Michelsohn,
\emph{Spin Geometry},
Princeton Mathematical Series 38,
Princeton University Press, Princeton, 1989.

\bibitem{LeeSmooth}
J.~M.~Lee,
\emph{Introduction to Smooth Manifolds},
2nd ed., Graduate Texts in Mathematics 218,
Springer, New York, 2013.

\bibitem{MilnorStasheff}
J.~W.~Milnor and J.~D.~Stasheff,
\emph{Characteristic Classes},
Annals of Mathematics Studies 76,
Princeton University Press, Princeton, 1974.

\bibitem{Nakahara}
M.~Nakahara,
\emph{Geometry, Topology and Physics},
2nd ed., CRC Press, Boca Raton, 2003.
ISBN 978-0-7503-0606-5.

\bibitem{ONeill}
B.~O'Neill,
\emph{Semi-Riemannian Geometry with Applications to Relativity},
Pure and Applied Mathematics 103,
Academic Press, New York, 1983.

\bibitem{Steenrod}
N.~Steenrod,
\emph{The Topology of Fibre Bundles},
Princeton Mathematical Series 14,
Princeton University Press, Princeton, 1951.

\bibitem{LeanSystem}
L.~de Moura, S.~Kong, J.~Avigad, F.~van Doorn, and J.~von Raumer,
The Lean theorem prover (system description),
in A.~P.~Felty and A.~Middeldorp (eds.),
\emph{Automated Deduction -- CADE-25},
Lecture Notes in Computer Science 9195,
Springer, Cham, 2015, pp.~378--388.

\bibitem{Mathlib2020}
The mathlib Community,
The Lean mathematical library,
in \emph{Proceedings of the 9th ACM SIGPLAN International Conference
on Certified Programs and Proofs (CPP 2020)},
ACM, 2020, pp.~367--381.
doi:10.1145/3372885.3373824.

\bibitem{RSCHFormal}
{\raggedright
\emph{The Relational Spin-Closure Hypothesis: formal Lean development},
source repository, commit
\path{e89832ec0e1f301c051b01a0a4118d9082118a60},
\url{https://github.com/antaris82/The_Relational_Spin-Closure_Hypothesis},
accessed 12 September 2026.\par}

\end{thebibliography}

\end{document}
